% ===== PRÉAMBULE CANONIQUE — à coller VERBATIM en tête de CHAQUE .tex =====
\documentclass[11pt,a4paper]{article}
\usepackage[utf8]{inputenc}\usepackage[T1]{fontenc}\usepackage{lmodern}
\usepackage[french]{babel}
\usepackage{amsmath,amssymb,mathtools}\usepackage{icomma}
\usepackage{siunitx}\sisetup{locale=FR}  % \qty{}{}, \si{}, \ang{}
\usepackage{xcolor}\usepackage{colortbl}
\usepackage{tikz}\usepackage{pgfplots}\pgfplotsset{compat=1.18}
\usepackage[siunitx,europeanresistors]{circuitikz} % circuits (élec.)
\usepackage[most]{tcolorbox}\usepackage{enumitem}\usepackage{array,booktabs}
\usepackage[a4paper,margin=2.2cm]{geometry}
\usetikzlibrary{arrows.meta,calc,angles,quotes,positioning,patterns,%
  backgrounds,shapes.geometric,decorations.pathmorphing,decorations.markings,intersections}
\definecolor{primary}{RGB}{222,49,99}\definecolor{primarydark}{RGB}{150,22,55}
\definecolor{defcol}{RGB}{0,114,178}\definecolor{methcol}{RGB}{52,92,148}
\definecolor{excol}{RGB}{0,135,90}\definecolor{attncol}{RGB}{200,40,55}
\tcbset{boxbase/.style={enhanced,breakable,boxrule=0.9pt,arc=2.5pt,
  left=3mm,right=3mm,top=2.3mm,bottom=2.3mm,fonttitle=\bfseries,coltitle=white}}
% --- boîtes TUTORIEL (noms EXACTS, ne pas renommer) ---
\newtcolorbox{cle}[1][]{boxbase,colback=defcol!6,colframe=defcol,
  title={\textsc{À retenir}\ifx\relax#1\relax\relax\else\textmd{ : \itshape #1}\fi}}
\newtcolorbox{methode}[1][]{boxbase,colback=methcol!7,colframe=methcol,sharp corners,
  title={\textsc{Méthode}\ifx\relax#1\relax\relax\else\textmd{ --- \itshape #1}\fi}}
\newtcolorbox{exemplebox}[1][]{boxbase,colback=excol!5,colframe=excol,
  title={\textsc{Exemple}\ifx\relax#1\relax\relax\else\textmd{ : \itshape #1}\fi}}
\newtcolorbox{piege}[1][]{boxbase,colback=attncol!6,colframe=attncol,
  title={\textsc{Piège}\ifx\relax#1\relax\relax\else\textmd{ : \itshape #1}\fi}}
\newtcolorbox{remarque}[1][]{boxbase,colback=black!4,colframe=black!45,
  title={\textsc{Remarque}\ifx\relax#1\relax\relax\else\textmd{ : \itshape #1}\fi}}
% --- boîtes EXERCICE / CORRIGÉ (noms EXACTS, auto-comptées) ---
\newtcolorbox[auto counter]{exo}[1][]{boxbase,colback=primary!4,colframe=primary,
  title={\textsc{Exercice}~\thetcbcounter\ifx\relax#1\relax\relax\else\textmd{ --- \itshape #1}\fi}}
\newtcolorbox[auto counter]{corrige}[1][]{boxbase,colback=excol!4,colframe=excol,
  title={\textsc{Corrigé}~\thetcbcounter\ifx\relax#1\relax\relax\else\textmd{ --- \itshape #1}\fi}}
\newcommand{\vv}[1]{\vec{#1}}\newcommand{\ds}{\displaystyle}
\newcommand{\R}{\mathbb{R}}\newcommand{\N}{\mathbb{N}}\newcommand{\Z}{\mathbb{Z}}
% (un \usepackage{fancyhdr}+en-tête est possible ; il est ignoré côté web)
% =========================================================================
\begin{document}
% THEME_TITLE: Le champ électrique
\sisetup{per-mode=symbol}

\begin{center}\large\bfseries\color{primarydark}
Thème 3 — Le champ électrique
\end{center}

Plutôt que de parler de \guillemotleft~force à distance~\guillemotright, on dit qu'une charge
\emph{modifie l'espace} autour d'elle : elle y crée un \textbf{champ électrique} $\vv{E}$. Une autre
charge placée dans ce champ subit alors la force $\vv{F} = q\,\vv{E}$.

\begin{cle}[définition du champ]
Le champ électrique en un point $M$ est la force par unité de charge subie par une charge test $q$ :
\[ \boxed{\;\vv{E} = \dfrac{\vv{F}}{q}\;}\qquad\Longleftrightarrow\qquad \vv{F} = q\,\vv{E}. \]
$\vv{E}$ \textbf{ne dépend pas} de la charge test : il ne dépend que des charges qui le créent et du
point $M$. Unité : $\si{\newton\per\coulomb} = \si{\volt\per\metre}$.
\end{cle}

\begin{cle}[champ d'une charge ponctuelle]
Une charge $Q$ crée à la distance $r$ un champ d'intensité
\[ \boxed{\,E = K\,\dfrac{|Q|}{r^{2}}\,},\qquad K \approx \num{9e9}\ \si{\newton\metre\squared\per\coulomb\squared}. \]
\textbf{Orientation} : $\vv{E}$ \emph{fuit} les charges positives (pointe vers l'extérieur) et
\emph{pointe vers} les charges négatives.
\end{cle}

\begin{center}
\begin{tikzpicture}[scale=1,>=Stealth,font=\footnotesize,
    ch/.style={circle,minimum size=8mm,inner sep=0pt,text=white,font=\bfseries\small}]
  % charge + : E fuit
  \node[ch,fill=attncol] at (-3.2,0) {$+$};
  \foreach \a in {0,45,...,315}{\draw[->,excol,line width=0.9pt] (-3.2,0)+(\a:0.55) -- +(\a:1.25);}
  \node[excol,font=\scriptsize] at (-3.2,-1.7) {$\vv{E}$ fuit le $+$};
  % charge - : E pointe vers
  \node[ch,fill=defcol] at (3.2,0) {$-$};
  \foreach \a in {0,45,...,315}{\draw[->,excol,line width=0.9pt] (3.2,0)+(\a:1.25) -- +(\a:0.55);}
  \node[excol,font=\scriptsize] at (3.2,-1.7) {$\vv{E}$ pointe vers le $-$};
\end{tikzpicture}
\end{center}

\begin{cle}[sens de la force selon le signe de $q$]
$\vv{F} = q\,\vv{E}$ : pour une charge test \textbf{positive}, $\vv{F}$ et $\vv{E}$ ont le
\emph{même} sens ; pour une charge \textbf{négative}, ils sont de sens \emph{opposés}.
\end{cle}

\begin{methode}[superposition et champ nul]
Le champ total est la \textbf{somme vectorielle} des champs de chaque charge : $\vv{E}_{\text{tot}} =
\sum_i \vv{E}_i$.
\begin{itemize}[leftmargin=1.5em,topsep=1pt,itemsep=1pt]
  \item \textbf{Champ nul} entre deux charges de \emph{même signe} : plus près de la petite charge,
        avec $\dfrac{|Q_1|}{x^2} = \dfrac{|Q_2|}{(L-x)^2}$ (prendre la racine).
  \item \textbf{Symétrie} : au centre d'un carré (ou d'un polygone régulier) de charges identiques,
        $\vv{E}_{\text{tot}} = \vv{0}$.
  \item Sur un \textbf{axe de symétrie} (médiatrice), les composantes transversales se compensent :
        il ne reste que la composante portée par l'axe.
\end{itemize}
\end{methode}

\begin{exemplebox}[champ d'une charge ponctuelle]
$Q = \SI{+5}{\nano\coulomb}$, $r = \SI{30}{\centi\metre}$ :
\[ E = \num{9e9}\times\frac{\num{5e-9}}{(0{,}30)^2} = \frac{45}{0{,}09} = \SI{500}{\newton\per\coulomb}, \]
dirigé \emph{vers l'extérieur} (la charge est positive).
\end{exemplebox}

\begin{piege}[ne pas confondre force et champ]
\begin{itemize}[leftmargin=1.5em,topsep=1pt,itemsep=1pt]
  \item La \textbf{force} dépend de la charge test ; le \textbf{champ} n'en dépend pas.
  \item $\vv{E} = \vv{F}/q$ et \emph{non} $E = Q\times F$ : le champ est la force \emph{divisée} par
        la charge.
  \item Un champ ne dépend jamais de la \textbf{masse} des particules, seulement des charges et des
        distances.
\end{itemize}
\end{piege}

\end{document}
